\chapter{Liver Function Tests (LFTs)}

\section*{What Is This Test?}
Liver function tests (LFTs), also called the liver panel or hepatic panel, are a group of blood tests that provide information about the functional and structural integrity of the liver. Despite the name, most components of the ``LFT'' are actually markers of hepatocyte injury or cholestasis, not direct measures of liver function. The standard LFT panel includes: (1) Hepatocellular injury markers: alanine aminotransferase (ALT) and aspartate aminotransferase (AST); (2) Cholestatic markers: alkaline phosphatase (ALP) and gamma-glutamyl transferase (GGT); (3) Excretory function marker: total and direct bilirubin (the liver's ability to conjugate and excrete bilirubin); (4) Synthetic function markers: albumin (hepatic protein synthesis) and prothrombin time/INR (synthesis of coagulation factors II, V, VII, IX, X, all vitamin K-dependent except factor V). Only albumin, INR, and bilirubin truly measure liver function---the aminotransferases and ALP measure hepatobiliary injury. The LFT is one of the most commonly ordered laboratory panels, typically performed as part of the comprehensive metabolic panel (CMP). Interpretation requires recognizing three patterns: hepatocellular (ALT/AST > ALP elevation), cholestatic (ALP/GGT > ALT/AST elevation), and mixed. Additionally, the severity of liver disease is assessed by synthetic function: decreasing albumin and rising INR indicate worsening hepatic reserve.

\section*{Alternative Names}
Hepatic Panel; Liver Panel; Liver Chemistries; Hepatic Function Panel; LFT; Liver Tests; Liver Enzymes; Comprehensive Metabolic Panel Liver Components

\section*{Principle}
The LFT is a panel of individual tests, each measured by different principles: ALT and AST are measured by kinetic spectrophotometry (IFCC methods, NADH consumption at 340 nm). ALP is measured by p-nitrophenylphosphate hydrolysis (absorbance at 405 nm). GGT is measured by gamma-glutamyl-3-carboxy-4-nitroanilide substrate hydrolysis (absorbance at 410 nm). Total bilirubin is measured by the diazo reaction (Jendrassik-Gr\'{o}f method, absorbance at 540--560 nm). Albumin is measured by dye-binding (bromocresol green or bromocresol purple, absorbance at 628 nm) or immunochemical methods. Prothrombin time/INR is a separate coagulation test performed on citrate plasma (blue-top tube), measuring the extrinsic and common coagulation pathways. The entire LFT panel (excluding INR, which requires a separate sample) is run simultaneously on a single automated chemistry analyzer platform (Roche Cobas, Abbott Architect, Siemens Atellica) from a single serum separator tube, with results available in 30--60 minutes \cite{tietz2023textbook}. This test panel is NOT performed by hematology analyzers.

\section*{History}
The modern LFT panel evolved from the work of multiple pioneers. The first liver function test was the measurement of bile acids and bilirubin in urine, used since ancient times (foam test, urobilinogen). Bromsulphthalein (BSP) clearance was introduced in the 1920s as a measure of hepatic excretory function. Serum albumin measurement became routine in the 1940s. The transaminases (AST, ALT) were discovered and applied clinically by Karmen, Wr\'{o}blewski, and LaDue in the 1950s. The De Ritis ratio was described in 1957. ALP, discovered in 1907, became part of the liver panel in the 1960s. GGT was added in the 1970s. The comprehensive metabolic panel standardized through CPT coding in the 1980s brought together all these tests. The MELD score, using bilirubin, INR, and creatinine, was developed in 2000 and adopted for liver transplant allocation in 2002 \cite{kamath2007model}.

\section*{Why This Test Is Ordered}
\begin{itemize}
  \item Routine screening for liver disease as part of annual physical examinations and preoperative assessment (CMP) \cite{tierney2023current}
  \item Evaluation of symptoms suggesting liver disease: jaundice, right upper quadrant pain, pruritus, dark urine, pale/clay-colored stools, nausea, fatigue, unexplained weight loss, spider angiomas, palmar erythema, gynecomastia, testicular atrophy, ascites, caput medusae
  \item Monitoring of known liver disease: viral hepatitis (B and C), alcoholic liver disease, NAFLD/NASH, autoimmune hepatitis, primary biliary cholangitis, primary sclerosing cholangitis, hemochromatosis, Wilson disease, alpha-1 antitrypsin deficiency
  \item Screening for hepatotoxicity from medications: acetaminophen, isoniazid, methotrexate, valproic acid, amiodarone, statins, anti-TNF agents, herbal and dietary supplements
  \item Assessment of disease severity in cirrhosis using synthetic function markers (albumin, INR, bilirubin) and calculation of prognostic scores: Child-Pugh score and MELD/MELD-Na score
  \item Evaluation of unexplained coagulopathy (elevated INR) not due to warfarin therapy
  \item Diagnosis of drug-induced liver injury (DILI): Hy's Law---ALT >3$\times$ ULN + bilirubin >2$\times$ ULN (in the absence of cholestasis, viral hepatitis, or other cause) indicates high risk of acute liver failure (10--50\% mortality) \cite{temple2006hys}
  \item Pre-transplant evaluation and post-transplant monitoring
\end{itemize}

\section*{Primary vs Secondary Test}
The LFT/CMP is a primary screening test performed on virtually all hospitalized patients and in most outpatient evaluations. It is the cornerstone of liver disease screening, diagnosis, and monitoring.

\section*{How This Test Is Performed}
A healthcare professional draws blood from a vein into a serum separator tube (gold-top) for the chemistry panel, plus a citrate (blue-top) tube if INR/PT is ordered. Standard venipuncture technique is used. For the chemistry panel, a single tube provides sufficient serum for all components (ALT, AST, ALP, GGT, bilirubin, albumin). The sample is centrifuged and analyzed on an automated chemistry analyzer. Results are available within 30--60 minutes. For the full LFT including INR, two tubes are needed. For bilirubin fractionation (direct and indirect), the same serum sample is used.

\section*{What Preparation Is Needed}
Fasting for 8--12 hours is recommended for accurate interpretation, especially for bilirubin (fasting can elevate unconjugated bilirubin in Gilbert syndrome) and GGT (fatty meals). Avoid alcohol for at least 24 hours before testing, as alcohol acutely elevates AST and GGT. Inform the provider of all medications, particularly: hepatotoxic drugs (acetaminophen, isoniazid, valproic acid, amiodarone, methotrexate, statins), enzyme-inducing drugs (phenytoin, phenobarbital, rifampin, carbamazepine---elevate GGT), drugs causing cholestasis (estrogens, anabolic steroids, erythromycin), anticoagulants (warfarin---elevates INR), and recent IV albumin infusion (falsely elevates measured albumin). Strenuous exercise within 24--48 hours can elevate AST, ALT, and CK from muscle release.

\section*{Contraindications and Precautions}
No absolute contraindications. Critical considerations for LFT interpretation: (1) LFTs are not truly ``liver function tests''---ALT, AST, ALP, and GGT measure injury, not function. Only albumin, INR, and bilirubin measure actual function \cite{kwo2017acg}. (2) A normal LFT does not exclude cirrhosis---patients with compensated cirrhosis may have normal LFTs. Always correlate with clinical examination (stigmata of chronic liver disease), platelet count (thrombocytopenia in portal hypertension from splenic sequestration), and imaging. (3) LFT patterns, not individual values, are most informative: hepatocellular pattern (ALT/AST > ALP), cholestatic pattern (ALP/GGT > ALT/AST), mixed pattern, and pattern of synthetic dysfunction (low albumin, high INR). (4) Non-hepatic causes of LFT abnormalities: muscle injury (AST > ALT), bone disease (ALP with normal GGT), hemolysis (elevated indirect bilirubin), malnutrition and nephrotic syndrome (low albumin), vitamin K deficiency and warfarin (elevated INR independent of liver function). (5) Macroenzymes (macro-AST, macro-ALP) can cause isolated, persistent elevations \cite{wallach2022interpretation}. (6) Pregnancy alters LFTs: ALP increases (placental), albumin decreases (hemodilution), AST/ALT may decrease slightly. Elevated AST/ALT in pregnancy suggests HELLP syndrome, acute fatty liver of pregnancy, or preeclampsia.

\section*{Risks}
Minimal risk from venipuncture. Mild pain or bruising resolves in 1--2 days.

\section*{What Happens After the Test}
Results available within 30--60 minutes. LFTs are interpreted according to the pattern of abnormality: (1) Hepatocellular pattern: ALT/AST elevation disproportionally greater than ALP. Causes: viral hepatitis, NAFLD/NASH, alcoholic hepatitis, drug-induced, autoimmune hepatitis, ischemic hepatitis, hemochromatosis, Wilson disease, alpha-1 antitrypsin deficiency \cite{giannini2005liver}. (2) Cholestatic pattern: ALP elevation disproportionality greater than ALT/AST. Causes: biliary obstruction (stone, stricture, malignancy), PBC, PSC, drug-induced cholestasis, infiltrative disease (sarcoid, lymphoma, TB), TPN, sepsis. (3) Mixed pattern: roughly equal elevations. (4) Isolated hyperbilirubinemia: elevated bilirubin with otherwise normal LFTs---suggests Gilbert syndrome (unconjugated), hemolysis (unconjugated), or Dubin-Johnson/Rotor syndromes (conjugated). (5) Synthetic dysfunction: low albumin, elevated INR, elevated bilirubin---indicates decompensated liver disease. Child-Pugh score classification uses bilirubin, albumin, INR, ascites, and encephalopathy to quantify cirrhosis severity (Class A: 5--6, Class B: 7--9, Class C: 10--15) \cite{pugh1973transection}. MELD score uses bilirubin, creatinine, INR, and sodium to prioritize liver transplant allocation \cite{kamath2007model}.

\section*{Understanding Results}

\subsection*{Below Normal}
\begin{itemize}
  \item \textbf{Low ALT/AST:} Not clinically significant. May be seen in elderly, chronic illness, uremia, end-stage liver disease (burned-out cirrhosis where hepatocyte mass is depleted).
  \item \textbf{Low ALP:} Hypophosphatasia, zinc deficiency, hypothyroidism, pernicious anemia, Wilson disease (especially fulminant presentation), or EDTA/citrate tube contamination.
  \item \textbf{Low bilirubin:} Normal. No disease associations.
  \item \textbf{Low albumin (<3.5 g/dL):} This IS clinically significant---indicative of decreased hepatic synthesis (cirrhosis, chronic liver disease), protein loss (nephrotic syndrome, protein-losing enteropathy), malnutrition/cachexia, or chronic inflammation (negative acute phase reactant). Albumin <2.8 g/dL: severe hepatic impairment.
  \item \textbf{Elevated INR (>1.2):} Indicates decreased hepatic synthesis of coagulation factors (cirrhosis, acute liver failure), vitamin K deficiency (malabsorption, prolonged antibiotics, warfarin therapy), or consumptive coagulopathy (DIC). INR >1.5 is a marker of liver dysfunction.
\end{itemize}

\subsection*{Above Normal}
\begin{itemize}
  \item \textbf{Hepatocellular pattern of injury:} See ALT and AST chapters for detailed differential diagnosis. Brief overview: Mild transaminitis (1--3$\times$ ULN): NAFLD, chronic hepatitis C, medications, alcohol. Moderate (3--15$\times$ ULN): acute viral hepatitis, autoimmune hepatitis, Wilson disease. Severe (>15$\times$ ULN): acute viral hepatitis, acetaminophen toxicity, ischemic hepatitis, autoimmune hepatitis.
  \item \textbf{Cholestatic pattern:} See ALP and GGT chapters. Biliary obstruction (choledocholithiasis, malignancy, PSC, PBC), drug-induced cholestasis, infiltrative disease, sepsis, TPN.
  \item \textbf{Elevated bilirubin:} See bilirubin chapter. Unconjugated: hemolysis, Gilbert, Crigler-Najjar. Conjugated: biliary obstruction, hepatocellular injury, Dubin-Johnson, Rotor, sepsis.
  \item \textbf{Hy's Law for drug-induced liver injury:} ALT >3$\times$ ULN + bilirubin >2$\times$ ULN (direct >35\%) in absence of cholestasis or other cause $\rightarrow$ high risk of acute liver failure \cite{temple2006hys}.
\end{itemize}

\section*{Normal Results}
\begin{itemize}
  \item \textbf{ALT:} Males: 10--40 IU/L; Females: 7--35 IU/L
  \item \textbf{AST:} Males: 10--40 IU/L; Females: 9--32 IU/L
  \item \textbf{ALP:} 44--147 IU/L (adults); higher in children/adolescents
  \item \textbf{GGT:} Males: 8--61 IU/L; Females: 5--36 IU/L
  \item \textbf{Total bilirubin:} 0.3--1.2 mg/dL
  \item \textbf{Direct bilirubin:} 0.0--0.3 mg/dL
  \item \textbf{Albumin:} 3.5--5.0 g/dL
  \item \textbf{Total protein:} 6.0--8.3 g/dL
  \item \textbf{INR (Prothrombin time):} 0.9--1.1 (not on warfarin)
  \item \textbf{Child-Pugh score components:} Bilirubin, albumin, INR, ascites, encephalopathy
  \item \textbf{MELD score components:} Bilirubin, creatinine, INR, sodium
\end{itemize}

\section*{Follow-up Tests}
\begin{itemize}
  \item \textbf{Complete LFT with INR and albumin:} For synthetic function assessment.
  \item \textbf{Abdominal ultrasound with Doppler:} First-line imaging for fatty liver, biliary dilation, cirrhosis, portal hypertension.
  \item \textbf{Viral hepatitis serologies:} HAV IgM, HBsAg, anti-HBc IgM, anti-HCV, HCV RNA.
  \item \textbf{Autoimmune markers:} ANA, anti-smooth muscle antibody, anti-LKM1, anti-mitochondrial antibody, IgG.
  \item \textbf{Metabolic and genetic tests:} Ferritin/transferrin saturation (hemochromatosis), ceruloplasmin/24h urine copper (Wilson), alpha-1 antitrypsin phenotype.
  \item \textbf{Acetaminophen level:} If toxicity suspected.
  \item \textbf{Liver biopsy:} Gold standard for staging and diagnosis.
  \item \textbf{FibroScan (transient elastography):} Non-invasive fibrosis staging.
  \item \textbf{CT or MRI of abdomen:} For masses, metastases, detailed liver and biliary anatomy.
  \item \textbf{Upper endoscopy:} Screen for esophageal varices if cirrhosis suspected.
\end{itemize}

\section*{References}
\bibliographystyle{unsrtnat}
\bibliography{references}

\clearpage
\section*{Regulatory Status and Authorization Date}
Regulatory status applies to the specific assay, instrument, manufacturer, and intended use rather than to the test name alone. No single approval date can be assigned to this test category. For a commercial assay, verify the current product in the relevant regulator's database: the U.S. Food and Drug Administration (FDA) clearance, approval, or authorization; India's Central Drugs Standard Control Organisation (CDSCO) licence or permission; the European Union In Vitro Diagnostic Regulation (EU IVDR) CE certificate issued through the applicable conformity-assessment route; or the United Kingdom Medicines and Healthcare products Regulatory Agency (MHRA) registration and UKCA/CE status. Laboratory-developed tests may not have an individual FDA, CDSCO, EU IVDR, or MHRA approval date.
\clearpage
